\clearpage
\phantomsection%
\addcontentsline{toc}{chapter}{ABSTRACT}
\thispagestyle{fancy}

\begin{center}
	\large \bfseries \MakeUppercase{Abstract}\\
	\bigskip
	
	\normalsize \normalfont \justifying \singlespacing
	The development of social media has driven the emergence of various forms of digital content, including memes that combine visual and textual elements. Among these, there are memes containing indications of self-harm that may potentially have negative impacts on users. The main challenge in identifying self-harm memes lies in their implicit nature, as meaning often arises from the interaction between image and text rather than from a single modality. This study aims to develop a multimodal-based binary classification model to detect self-harm memes by integrating Contrastive Language-Image Pre-training (CLIP)  as an image encoder and Efficiently Learning an Encoder that Classifies Token Replacements Accurately (ELECTRA) as a text encoder through an intermediate fusion scheme. The research data were obtained from a combination of primary and secondary sources, which were then annotated using an ensemble pseudo-labeling approach based on a multimodal model. The research process includes pre-processing, development of unimodal and multimodal models, and evaluation using confusion matrix-based metrics. The results show that the multimodal model achieved an F1-score of 96,1\%, which is higher than the unimodal image model at 91,5\% and the unimodal text model at 92,0\%. However, for ambiguous data, especially those containing sarcasm, dark humor, or conflicts between text and visuals, the model still faces difficulties in determining the correct label. Overall, these findings indicate that the multimodal approach is more effective in understanding context and detecting self-harm memes compared to the unimodal approach.\par

    \vspace{20pt}
	\raggedright \textbf{Keywords:} self-harm, meme, multimodal, CLIP, ELECTRA
	
	\vfill
	
\end{center}
\clearpage
